% Une ligne commentaire débute par le caractère « % »

\documentclass[a4paper]{article}

% Options possibles : 10pt, 11pt, 12pt (taille de la fonte)
%                     oneside, twoside (recto simple, recto-verso)
%                     draft, final (stade de développement)

\usepackage[utf8]{inputenc}   % LaTeX, comprends les accents !
\usepackage[T1]{fontenc}      % Police contenant les caractères français
\usepackage[francais]{babel}
\usepackage{fullpage}
\usepackage{multicol}
\usepackage{hyperref}
\hypersetup{
    colorlinks=true,
    linkcolor=blue,
    filecolor=magenta,      
    urlcolor=red
    }
% \urlstyle{same} % ça sert à rien ce truc
\usepackage{bookmark}
\usepackage{blindtext}



\usepackage{graphicx}  % pour inclure des images
\graphicspath{ {rapport/img/} }

%\pagestyle{headings}        % Pour mettre des entêtes avec les titres
                              % des sections en haut de page

 \title{  TP2 : Les capteurs\\         % Les paramètres du titre : titre, auteur, date
  Programmation mobile}
\author{Mohamad Satea Almallouhi - Tony Nguyen\\
  \emph{M1 Génie Logiciel}\\
  Faculté des Sciences\\
Université de Montpellier.}
\date{5 mars 2024}



\begin{document}
    \maketitle
    \begin{center}
    \includegraphics[width=\textwidth]{makima}
    \end{center}

    % \begin{abstract}     % Résumé du travail
    %   \emph{Description très succinte du problème et des différentes étapes de réalisation}
    % \end{abstract}
    \newpage
    %\dominitoc  % initializer les minitoc
    \tableofcontents

    \newpage
    \begin{multicols}{2}
        [
            Faire une vidéo, le rapport avec des screenchot des résultats et du code et enfin un read.md(instruction). En plus, pour le bonus, faire une belle application, des tests unitaires, utiliser Kotlin, faire le rapport en Latex.
        ]
        \section*{Introduction}
        \addcontentsline{toc}{section}{Introduction}
        \paragraph{}
            Dans ce TP, nous allons l'utilisation des capteurs intégré dans nos smartphones.
            
            Nous allons voir comment manipuler les différents types de capteur comme le GPS, la boussole, le gyroscope, etc ...
        \paragraph{}
            Les sections du rapport suit les exercices.
        \section*{Démonstration}
        \addcontentsline{toc}{section}{Démonstration}
            \paragraph{}
            En ligne sur Youtube, à l'adresse URL \url{https://youtu.be/nQUkpSUjJlY} une démonstration vidéo de notre travail.
        \section{Liste des capteurs}
            Lorsque l'on déploie une application sur un smarphone, nous ne pouvons pas être sûr des capteurs que l'on aura a notre disposition. Dans cette section, nous allons regarder tout les capteurs disponible sur un appareil donné.
        \section{Détection de la disponibilité des capteurs}  
        \section{Accéléromètre}
            \paragraph{}
                Nous avons fait une petite application qui change la couleur en fond en fonction de l'accélération du téléphone.
        \section{Direction}
        \section{Secouer un appareil}
        \section{Proximité}
        \section{Géolocalisation}
            \paragraph{}
                Dans de nombreuse application, il est souhaitable de connaitre la position de l'utilisateur. Pour ce faire, nous allons manipuler le gps.
        \section{Icone}
            \paragraph{}
                Nous avons choisit de modifier l'îcone de notre application. Nous allons maintenant expliqué comment nous cela a été fait.
    \end{multicols}
\end{document}